\section{General linear groups in tangent categories}
\label{section:general-linear-groups}
In this section, we introduce general linear groups as objects that classify \emph{Kleisli parametrized linear automorphisms}, which are those morphisms $f\colon\A\times M\to\T^n\A$ which are linear in their first argument and that satisfy a suitable notion of invertibility. After introducing this concept, we prove that, as the name suggests, general linear groups are in fact group objects. We then establish that general linear groups can be constructed entirely from LEC objects $\lnr_\A$ via a tangent triple equalizer.



%__________________________________________________________________________
\subsection{General linear group objects}
\label{subsection:general-linear-groups}
Corollary~\ref{corollary:internal-linear-homs-classify-kleisli-multilinear-morphisms} establishes that the LEC object $\lnr_\A$ classifies Kleisli parametrized linear morphisms $f\colon\A\times M\to\T\A$. We shall define general linear group objects as those objects that classify those parametrized linear morphisms $f\colon\A\times M\to\T\A$ which are \emph{invertible}. The issue is that we do not expect any $f\colon\A\times M\to\T\A$ to be invertible since $\A$ and $\T\A$ will not be isomorphic objects. To find the correct notion of invertibility, we shall interpret a morphism $f\colon\A\times M\to\T\A$ as an \emph{endomorphism} in the appropriate category.

\begin{definition}
\label{definition:parametrized-linear-isomorphism}
Given two differential objects $\A$ and $\B$, an \textbf{$M$-parametrized linear isomorphism} consists of a morphism $f\colon\A\times M\to\B$ that satisfies the following conditions:
\begin{description}
\item[PLI.1] $f$ is an $M$-parametrized linear morphism, in that:
\begin{equation*}
% https://q.uiver.app/#q=WzAsNCxbMCwwLCJcXFQoXFxBXFx0aW1lcyBNKSJdLFsxLDAsIlxcVFxcQiJdLFswLDEsIlxcQVxcdGltZXMgTSJdLFsxLDEsIlxcQiJdLFsyLDAsIlxcbGFtYmRhX1xcQVxcdGltZXMgel9NIl0sWzAsMSwiXFxUIGYiXSxbMywxLCJcXGxhbWJkYV9cXEIiLDJdLFsyLDMsImYiLDJdXQ==
\begin{tikzcd}
{\T(\A\times M)} & {\T\B} \\
{\A\times M} & \B
\arrow["{\T f}", from=1-1, to=1-2]
\arrow["{\lambda_\A\times z_M}", from=2-1, to=1-1]
\arrow["f"', from=2-1, to=2-2]
\arrow["{\lambda_\B}"', from=2-2, to=1-2]
\end{tikzcd}
\end{equation*}

\item[PLI.2] There exists a morphism $f^{-M}\colon\B\times M\to\A$, called \textbf{$M$-parametrized inverse}, that makes the following diagrams:
\begin{equation*}
% https://q.uiver.app/#q=WzAsNCxbMSwwLCJcXEFcXHRpbWVzIE1cXHRpbWVzIE0iXSxbMSwxLCJcXEJcXHRpbWVzIE0iXSxbMCwxLCJcXEFcXHRpbWVzIE0iXSxbMSwyLCJcXEEiXSxbMCwxLCJmXFx0aW1lc1xcaWRfTSJdLFsyLDAsIlxcaWRfXFxBXFx0aW1lc1xcRGVsdGFfTSJdLFsxLDMsImZeey1NfSJdLFsyLDMsIlxccGlfMSIsMl1d
\begin{tikzcd}
& {\A\times M\times M} \\
{\A\times M} & {\B\times M} \\
& \A
\arrow["{f\times\id_M}", from=1-2, to=2-2]
\arrow["{\id_\A\times\Delta_M}", from=2-1, to=1-2]
\arrow["{\pi_1}"', from=2-1, to=3-2]
\arrow["{f^{-M}}", from=2-2, to=3-2]
\end{tikzcd}\qquad
% https://q.uiver.app/#q=WzAsNCxbMCwwLCJcXEJcXHRpbWVzIE1cXHRpbWVzIE0iXSxbMCwxLCJcXEFcXHRpbWVzIE0iXSxbMSwxLCJcXEJcXHRpbWVzIE0iXSxbMCwyLCJcXEIiXSxbMiwwLCJcXGlkX1xcQlxcdGltZXNcXERlbHRhX00iLDJdLFsxLDMsImYiLDJdLFsyLDMsIlxccGlfMSJdLFswLDEsImZeey1NfVxcdGltZXNcXGlkX00iLDJdXQ==
\begin{tikzcd}
{\B\times M\times M} & \\
{\A\times M} & {\B\times M} \\
\B
\arrow["{f^{-M}\times\id_M}"', from=1-1, to=2-1]
\arrow["f"', from=2-1, to=3-1]
\arrow["{\id_\B\times\Delta_M}"', from=2-2, to=1-1]
\arrow["{\pi_1}", from=2-2, to=3-1]
\end{tikzcd}
\end{equation*}
commutative.
\end{description}
When $\A=\B$, we will refer to such an $f$ as a \textbf{parametrized linear automorphism} of $\A$.
\end{definition}

In Lemma~\ref{lemma:multilinearity-fibrationally}, we established that $M$-parametrized linear morphisms $f\colon\A\times M\to\B$ can be equivalently characterized as \emph{linear} morphisms in $\Simple_M[\X,\TT]$. The next lemma establishes a similar characterization for parametrized linear isomorphisms.

\begin{lemma}
\label{lemma:parametrized-linear-isomorphisms}
A morphism $f\colon\A\times M\to\B$ is a parametrized linear isomorphism if and only if it lifts to a linear isomorphism $f\colon\A\nto\B$ in $\Simple_M[\X,\TT]$. Furthermore, the $M$-parametrized inverse $f^{-M}$ of $f$ is unique.
\end{lemma}
\begin{proof}
From Lemma~\ref{lemma:multilinearity-fibrationally}, $f\colon\A\times M\to\B$ is an $M$-parametrized linear map if and only if $f\colon\A\nto\B$ is linear in $\Simple_M[\X,\TT]$. It is also easy to see that the conditions expressed by the diagrams of \textbf{[PLI.2]} establish that $f^{-M}$ is an inverse of $f$ in $\Simple_M[\X,\TT]$. Finally, since inverses of morphisms are always unique, it follows that parametrized inverses are also unique.
\end{proof}

For our goal, we need to impose an invertibility condition on morphisms of type $f\colon\A\times M\to\T\A$. However, if $f$ is an $M$-parametrized linear isomorphism, then $\A\cong\T\A$, which is not the case. Instead, we shall reinterpret this class of morphisms as endomorphisms in the Kleisli tangent category of $(\T,c)$.

\begin{definition}
\label{definition:kleisli-PLAs}
An \textbf{Kleisli $M$-parametrized linear isomorphism} of order $n$ consists of a morphism $f\colon\A\times M\to\T^n\B$ subject to the following conditions:
\begin{description}
\item[KPLI.1] $f$ is an $M$-parametrized linear morphism;

\item[KPLI.2] There exists a morphism $f^{-M}\colon\B\times M\to\T^n\A$, called \textbf{Kleisli $M$-parametrized inverse} of $f$, that satisfies the following conditions:
\begin{equation*}
% https://q.uiver.app/#q=WzAsNixbMCwwLCJcXEFcXHRpbWVzIE0iXSxbMSwwLCJcXEFcXHRpbWVzIE1cXHRpbWVzIE0iXSxbMiwwLCJcXFRebihcXEFcXHRpbWVzIE0pIl0sWzIsMSwiXFxUXm5cXFReblxcQSJdLFsyLDIsIlxcVF5uXFxBIl0sWzAsMiwiXFxBIl0sWzAsMSwiXFxpZFxcdGltZXNcXERlbHRhIl0sWzEsMiwiZlxcdGltZXMgel5uIl0sWzIsMywiXFxUXm5mXnstTX0iXSxbMyw0LCJcXG11Xm4iXSxbMCw1LCJcXHBpXzEiLDJdLFs1LDQsInpebiIsMl1d
\begin{tikzcd}
{\A\times M} & {\A\times M\times M} & {\T^n(\A\times M)} \\
&& {\T^n\T^n\A} \\
\A && {\T^n\A}
\arrow["{\id\times\Delta}", from=1-1, to=1-2]
\arrow["{\pi_1}"', from=1-1, to=3-1]
\arrow["{f\times z^n}", from=1-2, to=1-3]
\arrow["{\T^nf^{-M}}", from=1-3, to=2-3]
\arrow["{\mu^n}", from=2-3, to=3-3]
\arrow["{z^n}"', from=3-1, to=3-3]
\end{tikzcd}
\end{equation*}
\begin{equation*}
% https://q.uiver.app/#q=WzAsNixbMCwwLCJcXEFcXHRpbWVzIE0iXSxbMSwwLCJcXEFcXHRpbWVzIE1cXHRpbWVzIE0iXSxbMiwwLCJcXFRebihcXEFcXHRpbWVzIE0pIl0sWzIsMSwiXFxUXm5cXFReblxcQSJdLFsyLDIsIlxcVF5uXFxBIl0sWzAsMiwiXFxBIl0sWzAsMSwiXFxpZFxcdGltZXNcXERlbHRhIl0sWzEsMiwiZl57LU19XFx0aW1lcyB6Xm4iXSxbMiwzLCJcXFRebmYiXSxbMyw0LCJcXG11Xm4iXSxbMCw1LCJcXHBpXzEiLDJdLFs1LDQsInpebiIsMl1d
\begin{tikzcd}
{\A\times M} & {\A\times M\times M} & {\T^n(\A\times M)} \\
&& {\T^n\T^n\A} \\
\A && {\T^n\A}
\arrow["{\id\times\Delta}", from=1-1, to=1-2]
\arrow["{\pi_1}"', from=1-1, to=3-1]
\arrow["{f^{-M}\times z^n}", from=1-2, to=1-3]
\arrow["{\T^nf}", from=1-3, to=2-3]
\arrow["{\mu^n}", from=2-3, to=3-3]
\arrow["{z^n}"', from=3-1, to=3-3]
\end{tikzcd}
\end{equation*}
\end{description}
where $z^n$ and $\mu^n$ denote the unit and the multiplication of the colax tangent monad $(\T^n,c^n)$. When $\A=\B$, we will refer to such a map as a \textbf{Kleisli Parametrized Linear Automorphism}.
\end{definition}

\begin{lemma}
\label{lemma:kleisli-PLI}
An $M$-parametrized linear morphism $f\colon\A\times M\to\T^n\B$ is a Kleisli $M$-parametrized linear isomorphism if and only if $f$ lifts to a linear isomorphism $f\colon\A\nto\B$ in the Kleisli tangent category $\KL[\T_M^n,c^n_M]$ of the colax tangent monad $(\T_M^n,c_M^n)$, where $\T^n_M$ denotes the iterate of the tangent bundle functor $\T_M$ of $\Simple_M[\X,\TT]$. Moreover, the Kleisli parametrized inverse $f^{-M}$ of $f$ is unique.
\end{lemma}
\begin{proof}
That $f$ is a linear isomorphism in the Kleisli tangent category follows directly from the definition. Furthermore, uniqueness of the Kleisli parametrized inverse follows from the uniqueness of the inverse of an isomorphism.
\end{proof}

We can now define general linear group objects.

\begin{definition}
\label{definition:general-linear-group}
The \textbf{general linear group object} of $\A$ consists of an object $\GL_\A$ equipped with a map $\Ev\colon\A\times\GL_\A\to\A$, called the \textbf{evaluation map}, that satisfies the following conditions:
\begin{description}
\item[GL.1] $\Ev$ is a $\GL_\A$-parametrized linear automorphism.

\item[GL.2] $\Ev$ classifies Kleisli parametrized linear automorphisms of $\A$, in that if $f\colon\A\times M\to\T^n\A$ is a Kleisli $M$-parametrized linear automorphism of $\A$, there exists a unique $f\colon M\to\T^n\GL_\A$ that makes the following diagram
\begin{equation*}
% https://q.uiver.app/#q=WzAsMyxbMCwwLCJcXFRebihcXEFcXHRpbWVzXFxHTF9cXEEpIl0sWzEsMCwiXFxUXm5cXEEiXSxbMCwxLCJcXEFcXHRpbWVzIE0iXSxbMCwxLCJcXFReblxcRXYiXSxbMiwxLCJmIiwyXSxbMiwwLCJ6Xm5fXFxBXFx0aW1lc1xcaGF0IGYiXV0=
\begin{tikzcd}
{\T^n(\A\times\GL_\A)} & {\T^n\A} \\
{\A\times M}
\arrow["{\T^n\Ev}", from=1-1, to=1-2]
\arrow["{z^n\times\hat f}", from=2-1, to=1-1]
\arrow["f"', from=2-1, to=1-2]
\end{tikzcd}
\end{equation*}
commutative.
\end{description}
\end{definition}

\begin{example}
\label{example:GL-trivial}
Every category is trivially a tangent category whose tangent bundle functor and structural natural transformations are just the identity. When the category is Cartesian, it also becomes a Cartesian tangent category. However, the only differential object is the terminal object $\*$. It is also immediate to see that $\*$ always admits the general linear group, $\GL_\*$ being $\*$ and the evaluation map $\Ev_\*\colon\*\times\*\to\*$ being the unique terminal map.
\end{example}

\begin{example}
\label{example:GL-sman}
In the tangent category $\SMan$ of finitely-dimensional smooth manifolds, differential objects are precisely Euclidean spaces $\Real^n$ (\cite[Example~3.3(i)]{cockett:differential-bundles}). Every differential object admits the general linear group, whose $\GL$-object is precisely $\GL_n(\Real)$ and the evaluation map $\Real^n\times\GL_n(\Real)\to\Real^n$ sends each vector $v\in\Real^n$ and each linear automorphism $l\colon\Real^n\to\Real^n$ to $l(v)$.
\end{example}

\begin{example}
\label{example:GL-affine}
Affine schemes over a commutative and unital ring $R$ form a Cartesian tangent category $\Aff_R$ whose tangent bundle functor sends an affine scheme with coordinate ring $A$ to $\T^\Omega A\=\Sym_A\Omega_A$, that is, the symmetric algebra of the module of K\"ahler differentials of $A$. See~\cite{cruttwell:algebraic-geometry} for details. By~\cite[Corollary~4.18 and Remark~4.20]{cruttwell:algebraic-geometry}, there is an equivalence of categories:
\begin{align*}
&\DO[\Aff_R]\simeq\Mod_R^\op
\end{align*}
In particular, every differential object $\A$ of $\Aff_R$ is the free symmetric $R$-algebra $\A\cong\Sym_RM$ of an $R$-module $M$. For simplicity, we consider the case when $R$ is the field $\Real$ of real numbers. Consider a differential object $\A$ whose $\Real$-module $M$ is finitely generated by $n$ generators. Since $\Real$ is a field, $M$ becomes a vector space of dimension $n$.

Now, define $\GL_\A(-)$ as the functor that sends an $\Real$-algebra $A\in\Aff_\Real$ to the group $\GL_\A(A)$ of invertible $(n\times n)$-matrices with coefficients in $A$. This functor is represented by the $\Real$-algebra defined as follows:
\begin{align*}
&\GL_\A\=\Real[\{x_{i,j}\}_{i,j=1}^n][\det(x_{i,j})]^{-1}
\end{align*}
that is, the localization of the $\Real$-algebra freely-generated by the $n^2$-variables $x_{i,j}$, localized to invert the determinant polynomial $\det(x_{i,j})$ of the matrix $(x_{i,j})$. To see this, consider a morphism of $\Aff_\Real$ from $A\to\GL_\A$. This corresponds to a morphism of $\Real$-algebras $\GL_\A\to A$ which is fully determined by its values $a_{i,j}$ at each variable $x_{i,j}$. However, since it is an algebra homomorphism, $\det(a_{i,j})$ must be invertible in $A$. That is, the matrix $(a_{i,j})$ is invertible, thus, $(a_{i,j})\in\GL_\A(A)$.

The object $\GL_\A$ turns out to be the general linear group of $\A$.
\end{example}

To support our definition of general linear groups, we begin by showing that $\GL_\A$ classifies linear automorphisms of $\A$.

\begin{lemma}
\label{lemma:gl-classifies-linear-automorphisms}
The general linear group object $\GL_\A$ of a differential object $\A$ satisfies the following universal property:
\begin{description}
\item[GL.2'] For every $M$-parametrized linear automorphism $f\colon\A\times M\to\A$, there exists exactly one morphism $\hat f\colon M\to\GL_\A$ that makes the following diagram
\begin{equation*}
% https://q.uiver.app/#q=WzAsMyxbMCwwLCJcXEFcXHRpbWVzXFxHTF9cXEEiXSxbMSwwLCJcXEEiXSxbMCwxLCJcXEFcXHRpbWVzIE0iXSxbMCwxLCJcXEV2Il0sWzIsMSwiZiIsMl0sWzIsMCwiXFxpZF9cXEFcXHRpbWVzXFxoYXQgZiJdXQ==
\begin{tikzcd}
{\A\times\GL_\A} & \A \\
{\A\times M}
\arrow["\Ev", from=1-1, to=1-2]
\arrow["{\id_\A\times\hat f}", from=2-1, to=1-1]
\arrow["f"', from=2-1, to=1-2]
\end{tikzcd}
\end{equation*}
commutative.
\end{description}
\end{lemma}
\begin{proof}
Consider an $M$-parametrized linear automorphism $f\colon\A\times M\to\A$. We shall prove that $fz$ becomes a Kleisli $M$-parametrized linear automorphism of $\A$. For starters, since $z\T\lambda c=\lambda z_\T c=\lambda\T z$, it follows that $f$ is linear in the first argument. To prove that $fz$ is a Kleisli parametrized isomorphism, define $g\colon\A\times M\to\T\A$ as $f^{-M}z$, where $f^{-M}$ denotes the parametrized inverse of $f$. We shall compute:
\begin{align*}
(\id\times\Delta)(fz\times\id)(\id\times z)\T f^{-M}\T z\mu&=~(\id\times\Delta)(fz\times\id)(\id\times z)\T f^{-M}                      \tag{Unitality of the monad}\\
&=~(\id\times\Delta)(f\times\id)f^{-M}z             \tag{Naturality of $z$}\\
&=~(\id\times\Delta)\pi_1z                          \tag{$f^{-M}$ parametrized inverse of $f$}\\
&=~\pi_1z
\end{align*}
Similarly, we can also prove that $(\id\times\Delta)(f^{-M}z\times\id)(\id\times z)\T f\T z\mu=\pi_1z$. Therefore, $fz$ is a Kleisli $M$-parametrized linear automorphism of $\A$. Therefore, by the universal property \textbf{[GL.2]}, there exists a unique morphism $h\colon M\to\T\GL_\A$. Let $\hat f\colon M\to\GL_\A$, simply be $\hat f\=hp$. We shall prove that $\hat f$ satisfies $(\id\times\hat f)\Ev=f$:
\begin{align*}
(\id\times\hat f)\Ev&=~(zp\times hp)\Ev     \tag{Definition of $\hat f$ and $zp=\id$}\\
&=~(z\times h)\T\Ev p                       \tag{Naturality of $p$}\\
&=~fzp                                      \tag{Definition of $h$}\\
&=~f                                        \tag{$zp=\id$}
\end{align*}
Now, let us prove that $\hat f$ is unique. Let $t$ satisfy $(\id\times t)\Ev=f$. Thus:
\begin{align*}
&(z\times h)\T\Ev=fz=(\id\times t)\Ev z=(z\times tz)\T\Ev
\end{align*}
Thus, $tz=h$. Thus, $t=tzp=hp=\hat f$. This proves that $\hat f$ is unique.
\end{proof}

The universal property \textbf{[GL.2]} implies the universal property \textbf{[GL.2']} of Lemma~\ref{lemma:gl-classifies-linear-automorphisms}; however, it is strictly stronger, since the class of Kleisli parametrized automorphisms is bigger than the one of parametrized automorphisms. Although \textbf{[GL.2]} is strictly stronger than \textbf{[GL.2']}, we shall express \textbf{[GL.2]} entirely in terms of \textbf{[GL.2']}, using that Kleisli parametrized automorphisms are in fact Kleisli morphisms.

\begin{lemma}
\label{lemma:general-linear-groups-kleisli}
Consider a differential object $\A$ and an object $\GL_\A$ together with a map $\Ev\colon\A\times\GL_\A\to\A$. Thus, $\GL_\A$ is a general linear group for $\A$ with evaluation map $\Ev$ if and only if it satisfies the following conditions:
\begin{description}
\item[GL.1] $\Ev$ is a $\GL_\A$-parametrized linear automorphism of $\A$;

\item[GL.2''] For every $n>0$ and every parametrized linear automorphism $f\colon\A\times M\nto\A$ in $\KL[\T^n,c^n]$, there exists a unique morphism $\hat f\colon M\nto\GL_\A$ of $\KL[\T^n,c^n]$ that makes the following diagram
\begin{equation*}
% https://q.uiver.app/#q=WzAsMyxbMCwwLCJcXEFcXHRpbWVzXFxHTF9cXEEiXSxbMSwwLCJcXEEiXSxbMCwxLCJcXEFcXHRpbWVzIE0iXSxbMiwxLCJmIiwyLHsic3R5bGUiOnsiYm9keSI6eyJuYW1lIjoiYmFycmVkIn19fV0sWzAsMSwiXFxFdl5cXEtMIiwwLHsic3R5bGUiOnsiYm9keSI6eyJuYW1lIjoiYmFycmVkIn19fV0sWzIsMCwiXFxpZF9cXEFcXHRpbWVzXFxoYXQgZiIsMCx7InN0eWxlIjp7ImJvZHkiOnsibmFtZSI6ImJhcnJlZCJ9fX1dXQ==
\begin{tikzcd}
{\A\times\GL_\A} & \A \\
{\A\times M}
\arrow["{\Ev^\KL}"{inner sep=.8ex}, "\shortmid"{marking}, from=1-1, to=1-2]
\arrow["{\id_\A\times\hat f}"{inner sep=.8ex}, "\shortmid"{marking}, from=2-1, to=1-1]
\arrow["f"'{inner sep=.8ex}, "\shortmid"{marking}, from=2-1, to=1-2]
\end{tikzcd}
\end{equation*}
commutative in $\KL[\T^n,c^n]$, where $\SQ{\Ev^\KL}\=\Ev z^n$.
\end{description}
\end{lemma}
\begin{proof}
The argument for this statement is fairly similar to the one of~\ref{corollary:internal-linear-homs-classify-kleisli-multilinear-morphisms} as it consists entirely of rewriting condition \textbf{[GL.2]} as a condition in the Kleisli category of $\T^n$.
\end{proof}

By Lemma~\ref{lemma:gl-classifies-linear-automorphisms}, we can prove that $\GL_\A$ is a group object.

\begin{proposition}
\label{proposition:gl-group-object}
The general linear group object $\GL_\A$ of $\A$ is a group object whose unit, multiplication, and inverse maps are uniquely defined by the following commutative diagrams:
\begin{equation*}
% https://q.uiver.app/#q=WzAsMyxbMCwwLCJcXEFcXHRpbWVzXFxHTF9cXEEiXSxbMSwwLCJcXEEiXSxbMCwxLCJcXEFcXHRpbWVzXFwqIl0sWzAsMSwiXFxFdiJdLFsyLDEsIlxccGlfMSIsMl0sWzIsMCwiXFxpZF9cXEFcXHRpbWVzXFxlIl1d
\begin{tikzcd}
{\A\times\GL_\A} & \A \\
{\A\times\*}
\arrow["\Ev", from=1-1, to=1-2]
\arrow["{\id_\A\times\e}", from=2-1, to=1-1]
\arrow["{\pi_1}"', from=2-1, to=1-2]
\end{tikzcd}\qquad
% https://q.uiver.app/#q=WzAsNCxbMCwwLCJcXEFcXHRpbWVzXFxHTF9cXEEiXSxbMSwwLCJcXEEiXSxbMCwxLCJcXEFcXHRpbWVzXFxHTF9cXEFcXHRpbWVzXFxHTF9cXEEiXSxbMSwxLCJcXEFcXHRpbWVzXFxHTF9cXEEiXSxbMCwxLCJcXEV2Il0sWzIsMCwiXFxpZF9cXEFcXHRpbWVzXFxtIl0sWzMsMSwiXFxFdiIsMl0sWzIsMywiXFxFdlxcdGltZXNcXGlkIiwyXV0=
\begin{tikzcd}
{\A\times\GL_\A} & \A \\
{\A\times\GL_\A\times\GL_\A} & {\A\times\GL_\A}
\arrow["\Ev", from=1-1, to=1-2]
\arrow["{\id_\A\times\m}", from=2-1, to=1-1]
\arrow["{\Ev\times\id}"', from=2-1, to=2-2]
\arrow["\Ev"', from=2-2, to=1-2]
\end{tikzcd}\qquad
% https://q.uiver.app/#q=WzAsMyxbMCwwLCJcXEFcXHRpbWVzXFxHTF9cXEEiXSxbMSwwLCJcXEEiXSxbMCwxLCJcXEFcXHRpbWVzXFxHTF9cXEEiXSxbMCwxLCJcXEV2Il0sWzIsMSwiXFxFdl57LVxcR0xfXFxBfSIsMl0sWzIsMCwiXFxpZFxcdGltZXNcXGkiXV0=
\begin{tikzcd}
{\A\times\GL_\A} & \A \\
{\A\times\GL_\A}
\arrow["\Ev", from=1-1, to=1-2]
\arrow["{\id\times\i}", from=2-1, to=1-1]
\arrow["{\Ev^{-\GL_\A}}"', from=2-1, to=1-2]
\end{tikzcd}
\end{equation*}
where $\Ev^{-\GL_\A}$ denotes the parametrized inverse of $\Ev$.
\end{proposition}
\begin{proof}
One can directly check the equations that define a group object; however, we decide to give a more structural perspective on this result. We begin by constructing a functor $\GL_\A(-)\colon\X\to\Gr^\op$, where $\Gr$ is the category of groups. Let $\GL_\A(M)$ be the set of all $M$-parametrized linear automorphisms $f\colon\A\times M\to\A$. By Lemma~\ref{lemma:parametrized-linear-isomorphisms}, PLAs are linear isomorphisms in $\Simple_M[\X,\TT]$, thus, $\GL_\A(M)$ becomes a group object. Moreover, by Proposition~\ref{proposition:simple-tangent-fibration}, since $\Simple[\X,\TT]\to(\X,\TT)$ is a tangent fibration, the substitution functors $g^\*\colon\Simple_N[\X,\TT]\to\Simple_M[\X,\TT]$, induced by $g\colon M\to N$ are strong tangent morphisms. So, they preserve linear isomorphisms and, by functoriality, they restrict to group homomorphisms. Thus, we can define $\GL_\A(g)\=g^\*\colon\GL_\A(N)\to\GL_\A(M)$. It turns out that the universal property \textbf{[GL.2']}, established by Lemma~\ref{lemma:gl-classifies-linear-automorphisms} is equivalent to the representability of the functor $\GL_\A(-)$. We leave it to the reader to prove this statement. By the Yoneda lemma, the group structure of each $\GL_\A(M)$ induces a group structure on the representing object, which is precisely $\GL_\A$.
\end{proof}



%__________________________________________________________________________
\subsection{General linear groups as tangent equalizers}
\label{subsection:equalizers}
We now consider the case when $\A$ is a differential object which admits the LEC object $\lnr_\A$. We shall prove that one can construct the general linear group $\GL_\A$ of $\A$ out of $\lnr_\A$ via a \emph{tangent} triple equalizer. We begin with a technical lemma. Recall that a \textbf{tangent limit}, as introduced in~\cite[Definition~1.1.3]{macadam:vector-bundles}, consists of a limit cone $\gamma_c\colon L\to D(c)$ of a diagram $D\colon I\to\X$ which is preserved by all iterates of the tangent bundle functor, in that $\T^n\gamma_c\colon\T^nL\to\T D(c)$ is a limit cone for the diagram $\T^n\o D\colon I\to\X$. We shall also call a tangent limit of a particular shape, ``tangent [the name of that limit shape]'', e.g., a \textbf{tangent pullback} is a pullback diagram whose underlying cone is a tangent limit cone.

\begin{lemma}
\label{lemma:tangent-limits-kleisli}
Consider a diagram $D\colon I\to\X$ and a tangent structure $\TT$ on $\X$. A limit cone $\gamma_c\colon L\to D(c)$ for $D$ is a tangent limit if and only if, for every $n>0$, the inclusion functor $(\X,\TT)\to\KL[\T^n,c^n]$ preserves the limit cone $\gamma_c$.
\end{lemma}
\begin{proof}
This follows immediately from the fact that a limit cone $\gamma_c\colon L\to D(c)$ is preserved by a monad $S$ if and only if the inclusion functor $\X\to\KL[S]$ preserves such a limit cone.
\end{proof}

\begin{theorem}
\label{theorem:triple-equalizer}
Consider a differential object $\A$ and assume that the LEC object $\lnr_\A$ of $\A$ exists. Then, $\A$ admits the general linear group $\GL_\A$ if and only if the triple equalizer
\begin{equation*}
% https://q.uiver.app/#q=WzAsNCxbMCwwLCJcXEdMX1xcQSJdLFsxLDAsIlxcbG5yX1xcQVxcdGltZXNcXGxucl9cXEEiXSxbMywwLCJcXGxucl9cXEEiXSxbMiwwLCJcXCoiXSxbMSwyLCJcXG0iLDAseyJjdXJ2ZSI6LTR9XSxbMSwyLCJcXG1eXFxvcCIsMix7ImN1cnZlIjo0fV0sWzEsM10sWzMsMiwiXFxlIl0sWzAsMSwiIiwwLHsic3R5bGUiOnsiYm9keSI6eyJuYW1lIjoiZGFzaGVkIn19fV1d
\begin{tikzcd}
{\GL_\A} & {\lnr_\A\times\lnr_\A} & {\*} & {\lnr_\A}
\arrow[dashed, from=1-1, to=1-2]
\arrow[from=1-2, to=1-3]
\arrow["\m", curve={height=-24pt}, from=1-2, to=1-4]
\arrow["{\m^\op}"', curve={height=24pt}, from=1-2, to=1-4]
\arrow["\e", from=1-3, to=1-4]
\end{tikzcd}
\end{equation*}
exists and is a tangent triple equalizer.
\end{theorem}
\begin{proof}
We begin by showing that the triple equalizer exists if and only if there exists an object $\GL_\A$ together with a map $\Ev\colon\A\times\GL_\A\to\A$ that satisfies \textbf{[GL.1]} and \textbf{[GL.2']} of Lemma~\ref{lemma:gl-classifies-linear-automorphisms}. Using Lemma~\ref{lemma:general-linear-groups-kleisli}, we later prove that such a triple equalizer is a tangent triple equalizer if and only if $\GL_\A$ is a general linear group object with evaluation map $\Ev$.

For starters, let us assume that $\GL_\A$ and $\Ev$ satisfy \textbf{[GL.1]} and \textbf{[GL.2']} and let us prove that $\GL_\A$ is in fact a triple equalizer of $\m$, $!\e$, and $\m^\op$. To this end, we note that, since by \textbf{[GL.1]} $\Ev$ is a parametrized linear morphism, there exists a unique morphism $\theta\colon\GL_\A\to\lnr_\A$ that makes the following diagram
\begin{equation*}
% https://q.uiver.app/#q=WzAsMyxbMCwwLCJcXEFcXHRpbWVzXFxsbnJfXFxBIl0sWzEsMCwiXFxBIl0sWzAsMSwiXFxBXFx0aW1lc1xcR0xfXFxBIl0sWzIsMSwiXFxFdiIsMl0sWzAsMSwiXFxldiJdLFsyLDAsIlxcaWRfXFxBXFx0aW1lc1xcdGhldGEiXV0=
\begin{tikzcd}
{\A\times\lnr_\A} & \A \\
{\A\times\GL_\A}
\arrow["\ev", from=1-1, to=1-2]
\arrow["{\id_\A\times\theta}", from=2-1, to=1-1]
\arrow["\Ev"', from=2-1, to=1-2]
\end{tikzcd}
\end{equation*}
commutative. We shall leave it to the reader to prove that $\theta$ is in fact a monoid homomorphism. Thus, we shall define $\varphi\colon\GL_\A\to\lnr_\A\times\lnr_\A$ as $\varphi\=\<\theta,\i\theta\>$, where $\i\colon\GL_\A\to\GL_\A$ denotes the inverse map of the group object $\GL_\A$\footnote{Note that in Proposition~\ref{proposition:gl-group-object} we only used \textbf{[GL.2']} to prove that $\GL_\A$ is a group object.}. We shall prove that $\varphi$ equalizes $\m$, $!\e$, and $\m^\op$:
\begin{align*}
\varphi\m &=~\<\theta,\i\theta\>\m\\
&=~\<\id,\i\>\m\theta               \tag{$\theta$ preserves the multiplication}\\
&=~!\e\theta                        \tag{Inverse law of $\GL_\A$}\\
&=~!\e                              \tag{$\theta$ preserves the unit}
\end{align*}
Similarly, one can also prove that $\varphi\m^\op=!\e$. Now, consider a morphism $f=\<f_1,f_2\>\colon M\to\lnr_\A\times\lnr_\A$ that equalizes $\m$, $!\e$, and $\m^\op$. Consider also the two morphisms:
\begin{align*}
&\tilde f_k\colon\A\times M\xrightarrow{\id\times f_k}\A\times\lnr_\A\xrightarrow{\ev}\A,    &&\text{for }k=1,2.
\end{align*}
It is easy to see that $\tilde f_1$ and $\tilde f_2$ are parametrized linear endomorphisms of $\A$; we shall also prove that they are parametrized inverses to each other:
\begin{align*}
(\id\times\Delta)(\tilde f_1\times\id)\tilde f_2&=~(\id\times\<f_1,f_2\>)(\ev\times\id)\ev   \tag{Definitions of $\tilde f_1$ and $\tilde f_2$}\\
&=~(\id\times f)(\id\times\m)\ev                                    \tag{Definition of $f$ and of $\m$}\\
&=~(\id\times!\e)\ev                                                \tag{$f\m=!\e$}\\
&=~\pi_1                                                            \tag{Definition of $\e$}
\end{align*}
Similarly, using that $f\m^\op=!\e$, one can also show that $(\id\times\Delta)(\tilde f_2\times\id)\tilde f_1=\pi_1$, proving that $\tilde f_1$ is a parametrized linear automorphism of $\A$. Therefore, by \textbf{[GL.2']}, there exists a unique morphism $\hat f\colon M\to\GL_\A$ such that $(\id\times\hat f)\Ev=\tilde f_1$. However, we can also compute that:
\begin{align*}
&(\id\times f_1)\ev=\tilde f_1=(\id\times\hat f)\Ev=(\id\times(\hat f\theta)\ev
\end{align*}
Thus, by \textbf{[GL.2']}, $\hat f\theta=f_1$. Moreover, there exists also a unique morphism $\hat f^{-M}\colon M\to\GL_\A$ such that $(\id\times\hat f^{-M})\Ev=\tilde f_2$, since also $\tilde f_2$ becomes a parametrized linear automorphism of $\A$. However:
\begin{align*}
&(\id\times f_2)\ev=\tilde f_2=(\id\times\hat f^{-M})\Ev= (\id\times(\hat f^{-M}\theta))\ev
\end{align*}
Thus, $f_2=\hat f^{-M}\theta$. However, using the inverse law, $(\id\times(\hat f\i)\theta)\ev$ is a parametrized inverse for $\tilde f_1$. Thus, by uniqueness of the parametrized inverse, $(\id\times(\hat f\i)\theta)\ev=\tilde f_2$; thus, $f_2=\hat f^{-M}\theta=(\hat f\i)\theta$. Therefore, $\hat f\varphi=\hat f\<\theta,\i\theta\>=\<f_1,f_2\>=f$. Now, let us consider another morphism $g\colon M\to\GL_\A$ that satisfies $g\varphi=f$. Thus, $g\theta=f_1=\hat f\theta$. Thus, $(\id\times g)\Ev=(\id\times g\theta)\ev=(\id\times\hat f\theta)\ev=(\id\times\hat f)\Ev$. Thus, by \textbf{[GL.2']}, $g=\hat f$. This proves that $\GL_\A$ is the triple equalizer of $\m$, $!\e$, and $\m^\op$.

We shall now prove the converse; so, we shall assume that $\GL_\A$ is an object that comes with a map $\varphi\=\<\varphi_1,\varphi_2\>\colon\GL_\A\to\lnr_\A\times\lnr_\A$ that makes it into the triple equalizer of $\m$, $!\e$, and $\m^\op$. We shall prove that, under this assumption, $\GL_\A$ satisfies \textbf{[GL.1]} and \textbf{[GL.2']}. We begin by constructing the evaluation map $\Ev$ as follows:
\begin{align*}
&\Ev\colon\A\times\GL_\A\xrightarrow{\id\times\varphi_1}\A\times\lnr_\A\xrightarrow{\ev}\A
\end{align*}
Let also:
\begin{align*}
&\Ev^{-\GL_\A}\colon\A\times\GL_\A\xrightarrow{\id\times\varphi_2}\A\times\lnr_\A\xrightarrow{\ev}\A
\end{align*}
We already proved before that, if $f=\<f_1,f_2\>$ equalizes $\m$, $!\e$, and $\m^\op$, then $(\id\times f_1)\ev$ and $(\id\times f_2)\ev$ are in fact parametrized linear automorphisms and parametrized inverses of each other. Therefore, by construction, so are $\Ev$ and $\Ev^{-\GL_\A}$. Thus, $\Ev$ satisfies \textbf{[GL.1]}. To prove \textbf{[GL.2']}, consider a parametrized linear automorphism $f\colon\A\times M\to\A$ whose parametrized inverse is denoted by $f^{-M}$. Using the universal property \textbf{[ILH.2]}, we find two morphisms $\hat f,\hat f^{-M}\colon M\to\lnr_\A$ satisfying $(\id\times\hat f)\ev=f$ and $(\id\times\hat f^{-M})\ev$. We shall prove that $\<\hat f,\hat f^{-M}\>$ equalizes $\m$, $!\e$, and $\m^\op$:
\begin{align*}
(\id\times(\<\hat f,\hat f^{-M}\>\m)\ev&=~(\id\times\<\hat f,\hat f^{-M}\>)(\ev\times\id)\ev    \tag{Definition of $\m$}\\
&=~(\id\times\Delta)((\id\times\hat f)\ev\times\id)(\id\times\hat f^{-M})\ev\\
&=~(\id\times\Delta)(f\times\id)f^{-M}              \tag{Definitions of $\hat f$ and $\hat f^{-M}$}\\
&=~\pi_1                                            \tag{$f^{-M}$ parametrized inverse of $f$}\\
&=~(\id\times!\e)\ev                                \tag{Definition of $\e$}
\end{align*}
Thus, by \textbf{[ILH.2]}, $\<\hat f,\hat f^{-M}\>\m=!\e$. Similarly, one can also show that $\<\hat f,\hat f^{-M}\>\m^\op=!\e$. Thus, by the universal property of the triple equalizer, there exists a unique morphism $\tilde f\colon M\to\GL_\A$ such that $\tilde f\varphi=\<\hat f,\hat f^{-M}\>$. Thus, $\tilde\varphi_1=\hat f$. Therefore, $(\id\times\tilde f)\Ev=(\id\times\tilde f\varphi_1)\ev=(\id\times\hat f)\ev=f$. We leave it to the reader to prove that, in fact, $\tilde f$ is the unique morphism that satisfies this equation.

So far, we proved that $\GL_\A$ is the triple equalizer of $\m$, $!\e$, and $\m^\op$ if and only if it comes with a map $\Ev$ that satisfies \textbf{[GL.1]} and \textbf{[GL.2']}. To prove the theorem, we shall make use of Lemma~\ref{lemma:general-linear-groups-kleisli}. According to this lemma, $\GL_\A$ is a general linear group with evaluation map given by $\Ev$ if and only if $\Ev$ satisfies \textbf{[GL.1]} and it satisfies \textbf{[GL.2']} in each Kleisli category $\KL[\T^n,c^n]$. Therefore, $\GL_\A$ is a general linear group if and only if, in each Kleisli category $\KL[\T^n,c^n]$, $\GL_\A$ is the triple equalizer of $\m$, $!\e$, and $\m^\op$. However, by Lemma~\ref{lemma:tangent-limits-kleisli}, a limit cone is a tangent limit if and only if it is a limit cone in each Kleisli category $\KL[\T^n,c^n]$. Therefore, $\GL_\A$ is a general linear group if and only if $\GL_\A$ is the \emph{tangent} triple equalizer of $\m$, $!\e$, and $\m^\op$.
\end{proof}

Conceptually, Theorem~\ref{theorem:triple-equalizer} establishes that the general linear group $\GL_\A$ is the maximal subgroup of $\lnr_\A$ and that the embedding is compatible with the tangent bundle functor.



%__________________________________________________________________________
%__________________________________________________________________________